\chapter{Starting and stopping ProPrep}\label{ch:starting}

This chapter shows you how to open ProPrep, answer the questions it asks the
first time, close it safely, and come back to a project another day. You can
run ProPrep in two ways: inside a web browser, or in the Terminal. The screens
and questions are the same in both.

\section{Two ways to run ProPrep}\label{sec:starting-twoways}

\begin{tabularx}{\linewidth}{@{}lL@{}}
\toprule
\textbf{Command} & \textbf{What you get} \\
\midrule
\typed{proprep-web} & ProPrep in a browser tab, with the 3D structure viewer
docked beside it and a \textsf{Download project} button. Recommended for most
users. \\
\typed{proprep} & ProPrep directly in the Terminal window. The 3D viewer opens
in a separate browser tab when you ask for it. \\
\bottomrule
\end{tabularx}

\medskip
Both commands start the same program. Everything in this manual works the same
way in either one.

\subsection*{Before you start: choose a work folder}

ProPrep saves its files in the folder you start it from (or in a new folder
inside it). So decide first where your projects should live. A good choice is a
folder used only for ProPrep work, for example \file{ProPrep-work} in your home
folder.

Open the \term{Terminal} (the program where you type commands; Chapter~\ref{ch:installation}
explains how to open it on your computer). Then make the folder once and move
into it. The command \typed{cd} (``change directory'') moves you into a folder;
\typed{mkdir} (``make directory'') creates one.

\begin{terminal}
mkdir ~/ProPrep-work
cd ~/ProPrep-work
\end{terminal}

Here \file{~} is a short way of writing your home folder. Next time you only
need the \typed{cd} line. If you installed ProPrep with the conda script, also
type \typed{conda activate ProPrep} first (Chapter~\ref{ch:installation}).

\section{Starting in the browser (\texttt{proprep-web})}\label{sec:starting-web}

\begin{steps}
  \item In the Terminal, inside your work folder, type:
\begin{terminal}
proprep-web
\end{terminal}
  \item ProPrep prints a few lines like these, and then opens your web browser:
\begin{screen}
ProPrep Web Shell listening on http://127.0.0.1:8000/
  - Terminal pane runs the unmodified ProPrep CLI
  - Structure Viewer will dock on the right when launched from the CLI
  - Session transcript: ~/.proprep/web_sessions/session_<timestamp>.log
...
INFO:     Uvicorn running on http://127.0.0.1:8000 (Press CTRL+C to quit)
\end{screen}
  The address on the first line is where ProPrep is running. \texttt{127.0.0.1}
  means ``this computer'': nothing is sent over the internet. If the browser
  does not open by itself, copy the address into your browser.
  \item \textbf{Leave the Terminal window open.} It is running ProPrep for the
  browser tab. You can ignore the \texttt{INFO} lines that keep appearing there.
  \item Work in the browser tab. The left part of the page is the ProPrep
  terminal, where you type your answers. The right part is kept for the
  structure viewer, which appears there automatically when ProPrep shows a
  structure. The top bar has \textsf{Download project}, buttons to make the text
  bigger or smaller (\textsf{A$-$} and \textsf{A+}), and \textsf{Detach viewer}.
\end{steps}

\begin{tip}
If another program is already using port 8000, ProPrep takes the next free
number and says so, for example \texttt{(port 8420 was busy; using 8421)}.
Always read the address from the \texttt{listening on} line.
\end{tip}

The browser page loads part of its terminal display from the internet, so the
browser needs an internet connection even though ProPrep itself runs on your
computer.

\subsection*{Downloading your project}

\textsf{Download project} gives you a \file{.zip} file (named like
\file{proprep_local_20261121_125705.zip}) of your work. It contains
\emph{everything} in the folder you started \typed{proprep-web} from,
including all folders inside it. This is another reason to start
\typed{proprep-web} from a dedicated work folder and not from your home folder.
The files are also already on your computer, in that folder; the download is a
convenient copy.

\subsection*{Stopping the browser version}

The browser version stops by itself about two seconds after you close the last
ProPrep tab. The Terminal then shows
\texttt{No browser reconnected within 2.0s; shutting down.} and returns to its
normal prompt. You can also stop it by pressing \key{Ctrl}+\key{C} in the
Terminal window.

\begin{watchout}
Closing the tab ends ProPrep at once, without the ``Are you sure?'' question.
Your files and the session record are kept, but it is tidier to exit ProPrep
first with \key{x} (Section~\ref{sec:starting-exit}), use
\textsf{Download project} if you want a copy, and only then close the tab.
\end{watchout}

\begin{behindscenes}
Each browser session is also written to a plain-text log in
\file{~/.proprep/web_sessions/}. Start with \typed{proprep-web --no-transcript}
if you do not want this log. A few other options (colors, font size, port) are
listed in Chapter~\ref{ch:cmdline}.
\end{behindscenes}

A few things differ slightly in the browser. The structure viewer docks in the
page instead of opening a new tab, and the PDB Filter does not ask whether to
open the viewer; it shows the structure in the docked viewer instead. A session
recorded in the browser is marked as such, which matters only for batch
processing (Chapter~\ref{ch:batch}).

\section{Starting in the Terminal (\texttt{proprep})}\label{sec:starting-terminal}

To follow this manual, start ProPrep in \term{Full Menu} mode (the menu that
shows every tool at once; Section~\ref{sec:starting-modes}):

\begin{terminal}
proprep --full-menu
\end{terminal}

Typing \typed{proprep} on its own also works; ProPrep then asks which menu you
want (Section~\ref{sec:starting-questions}). The part after the command, such as
\typed{--full-menu}, is called an \term{option} or flag: an extra instruction
for this run only. Chapter~\ref{ch:cmdline} lists the useful ones.

\subsection*{The start-up messages}

Every start begins with a few lines like these:

\begin{screen}
MODELLER not found. Structure repair (missing loops/residues) will be unavailable.
  To install: conda install -c salilab modeller
  A free academic license is required: https://salilab.org/modeller/registration.html
  After registering, set your license key:
    export KEY_MODELLER='your_license_key_here'
  Or edit the config.py file shown during conda installation.
Successfully imported core modules
\end{screen}

The MODELLER message appears on every start when MODELLER has no license key. It
affects only the repair of missing loops and residues (Chapter~\ref{ch:cleaning}).
You can ignore it otherwise. Chapter~\ref{ch:installation} explains how to add
the key. Then the ProPrep banner appears, with the version number
(\texttt{v1.23.0}).

\section{The first-launch questions}\label{sec:starting-questions}

\subsection*{Which menu?}

Unless you started with \typed{--full-menu} (or \typed{--workflow}), ProPrep
shows a box called ``Choose Your Experience'' and asks:

\begin{screen}
 1. Workflow Mode (Recommended for new users)
 ...
 2. Full Menu Mode (For experienced users)
 ...
Select menu mode  [1/2]: §ans«2»

Save this as your default preference?  [y/n] (n): §ans«n»
Tip: Use --workflow or --full-menu flags to skip this prompt
\end{screen}

\begin{steps}
  \item At \prompt{Select menu mode [1/2]}, type \key{2} for Full Menu, which
  this manual uses. This question has \emph{no} default: pressing \key{Enter}
  alone only repeats it.
  \item At \prompt{Save this as your default preference?}, type \key{n} (or
  press \key{Enter}). If you answer \key{y}, ProPrep stops asking this question
  on later starts, and changing your mind means going into Preferences
  (Section~\ref{sec:menus-prefs}). Answering \key{n} keeps the choice open.
\end{steps}

\begin{tip}
Once you know you always want the Full Menu, you can save it as the default
with one command: \typed{proprep --set-default full-menu}. This is handy for
the browser version, which cannot be given \typed{--full-menu}.
\end{tip}

\subsection*{The project name}

Next comes ``Project Setup'':

\begin{screen}
═══ Project Setup ═══

Enter a name for your project. All structures and outputs
will be saved in a directory with this name.
Press Enter to use the current directory.

Project name (.): §ans«myproject»
✓ Project directory: /Users/.../ProPrep-work/myproject
\end{screen}

Every file ProPrep makes for this project goes into this folder.

\begin{itemize}
  \item \textbf{Type a name}, such as \typed{ubiquitin} or \typed{1ubq\_md}.
  ProPrep creates a folder with that name inside the folder you started from,
  and works inside it. If the folder already exists, ProPrep uses it and keeps
  the files already there.
  \item \textbf{Press \key{Enter}} to use the folder you started from (the
  \texttt{.} in \prompt{(.)} means ``this folder''). Do this when you started
  ProPrep from inside an existing project folder.
\end{itemize}

Rules for a good project name:
\begin{itemize}
  \item Use letters, numbers, \texttt{-} and \texttt{\_} only.
  \item Avoid spaces. ProPrep accepts \typed{my project}, but some Amber
  programs used later can have trouble with spaces in folder names.
  \item Never use \texttt{/}. A name such as \typed{sub/deeper} stops ProPrep
  with \texttt{Error: [Errno 2] No such file or directory: 'sub/deeper'}.
  ProPrep can create only one new folder level here.
\end{itemize}

\begin{watchout}
The menu choice and the project name are given before ProPrep starts recording
your answers. They are not in the session file, and \key{undo} cannot go back
to them. If you type \typed{undo} at \prompt{Project name}, ProPrep makes a
folder called \file{undo}.
\end{watchout}

\section{What ProPrep prints next}\label{sec:starting-next}

After the project folder is set, ProPrep reports that it is recording:

\begin{screen}
Session recording to: ./proprep_session_20261121_123119.json
Tip: Use --no-session to disable automatic session recording
Tip: Type 'undo' at any prompt to rewind to an earlier answer
\end{screen}

The \term{session file} (\file{proprep_session_<date>_<time>.json}) is
ProPrep's record of every answer you type from now on, saved inside the project
folder. It is the real record of your work: it lets you go back with
\key{undo}, resume later, and repeat the same preparation on other proteins
(Chapters~\ref{ch:sessions} and~\ref{ch:batch}). The file is written as you go,
so it survives even if ProPrep is closed suddenly.

\subsection*{The update check}

On the very first start, under the banner, ProPrep explains:

\begin{screen}
ProPrep checks GitHub once a day for a newer release. To turn this off:
Preferences (p) > Check for new ProPrep releases at startup.
\end{screen}

ProPrep asks the ProPrep download page at most once a day whether a newer
version exists. Nothing is downloaded or changed, and nothing about you is sent.
With no internet connection, nothing is printed. When a newer version exists,
ProPrep prints a bold line such as
\texttt{ProPrep X.Y.Z is available (you have 1.23.0).} followed by the update
command (we could not see this, because 1.23.0 was the newest version).
Chapter~\ref{ch:installation} explains updating. Copies of ProPrep that came
bundled with AmberTools never check.

Finally, the main menu appears. Chapter~\ref{ch:menus} explains how to read it.

\section{Workflow Mode or Full Menu?}\label{sec:starting-modes}

ProPrep has two kinds of main menu. Both open the same tools.

\begin{tabularx}{\linewidth}{@{}p{0.2\linewidth}LL@{}}
\toprule
 & \textbf{Workflow Mode} & \textbf{Full Menu} \\
\midrule
Shows & Only the tools of the current stage (Load, Detect, Compare, Fix,
Simulate), with a suggestion of what to do. & Every tool at once, grouped in
five sections. \\
Tool keys & Numbers \key{1}, \key{2}, \ldots{} that restart at every stage,
so \key{1} means a different tool at each stage. & Fixed codes such as
\key{1a} or \key{4e} that never change. \\
Capital letters & Accepted. & Not accepted (type lowercase). \\
Switch to the other & \key{a} (all tools) & \key{g} \\
Start directly & \typed{proprep --workflow} & \typed{proprep --full-menu} \\
\bottomrule
\end{tabularx}

\medskip
\textbf{This manual uses Full Menu keys}, because a code like \key{4e} always
means the same tool, whatever you did before. In Workflow Mode, \key{4e}
(Protonation State Analyzer) is \key{5} on the Fix stage, and you must first
reach that stage. If you prefer Workflow Mode, Section~\ref{sec:menus-workflow}
explains it, and Appendix~\ref{app:menumap} lists where each tool sits.

Switching with \key{a} or \key{g} lasts only for the current run. To change the
mode ProPrep starts in, use Preferences (Section~\ref{sec:menus-prefs}) or
\typed{proprep --set-default full-menu}.

\section{Exiting ProPrep}\label{sec:starting-exit}

\begin{steps}
  \item At any menu, type \key{x}.
  \item At \prompt{Are you sure you want to exit? [y/n] (n)}, type \key{y}.
  Pressing \key{Enter} means \key{n}: you stay in ProPrep.
\end{steps}

ProPrep says goodbye, stops recording, and shows a thank-you box:

\begin{screen}
Are you sure you want to exit? [y/n] (n): §ans«y»
Goodbye!
[Session recording stopped]
...
╭───────────────────── Thank you for using ProPrep ──────────────────────╮
│ If you found ProPrep useful, please cite:                              │
│                                                                        │
│ Walker, A. & Guberman-Pfeffer, M.J. (2026). ProPrep: An Interactive    │
│ and Instructional Interface for Proper Protein Preparation with AMBER. │
│ bioRxiv. https://doi.org/10.64898/2026.02.26.708365                    │
│                                                                        │
│ Feedback (positive or negative) is welcome at: mjgp@nmsu.edu           │
╰────────────────────────────────────────────────────────────────────────╯
\end{screen}

Your files stay in the project folder. In the browser version, you can now
close the tab.

\begin{itemize}
  \item \key{x} works inside a tool's menu too. If you then answer \key{n},
  ProPrep takes you to the \emph{main} menu, not back to the tool.
  \item Pressing \key{Ctrl}+\key{C} stops ProPrep at any question. It prints
  \texttt{Application interrupted by user} and the same thank-you box. The
  session file keeps everything up to that moment.
\end{itemize}

\section{Coming back to a project later}\label{sec:starting-resume}

Your files stay in the project folder after you exit. ProPrep's memory of what
it loaded and made, however, starts empty at every start. There are two ways to
carry on.

\subsection*{Resume the recorded session}

ProPrep looks for session files only in the folder you start it from. So to be
offered a resume, first move \emph{into} the project folder:

\begin{terminal}
cd ~/ProPrep-work/myproject
proprep --full-menu
\end{terminal}

(For the browser, type \typed{proprep-web} at this point instead.) Before the
project question, ProPrep then asks:

\begin{screen}
Previous session files found!

What would you like to do?
  1. Start fresh session (record new interactions)
  2. Resume a previous session from any point
  3. Continue without session recording

Select option  [1/2/3] (1): §ans«2»
                          Available Session Files
┏━━━━━┳━━━━━━━━━━━━━━━━━━━━━━━━━━━━━━━━━━━━━━┳━━━━━━━━━━━━━━━━━━┳━━━━━━━━━━━━━━┓
┃ #   ┃ Filename                             ┃ Date/Time        ┃ Interactions ┃
┡━━━━━╇━━━━━━━━━━━━━━━━━━━━━━━━━━━━━━━━━━━━━━╇━━━━━━━━━━━━━━━━━━╇━━━━━━━━━━━━━━┩
│ 1   │ proprep_session_20261121_123119.json │ 2026-11-21 12:32 │ 74           │
└─────┴──────────────────────────────────────┴──────────────────┴──────────────┘

Select session to resume (0 to cancel)  [1/0] (0): §ans«1»
\end{screen}

\begin{steps}
  \item Type \key{2} to resume.
  \item Type the number of the session you want. \textbf{Do not just press
  \key{Enter}}: the default here is \key{0}, which cancels and starts a fresh
  session instead.
  \item ProPrep then offers \prompt{1. Continue from end}, \prompt{2. Open
  session editor} and \prompt{3. Cancel}. Chapter~\ref{ch:sessions} explains
  each one.
\end{steps}

Resuming replays your recorded answers one by one to rebuild ProPrep's memory.
It repeats each step (downloads, filtering and so on), so it takes a little
time.

\begin{watchout}
In version 1.23.0, \prompt{Continue from end} on a session that you ended with
\key{x} also replays your exit, so ProPrep quits again straight away. Use the
session editor first to delete the last two answers (the \key{x} and the
\key{y}); Chapter~\ref{ch:sessions} shows how. Sessions that ended with a closed
browser tab or \key{Ctrl}+\key{C} have no recorded exit and continue normally.
\end{watchout}

\subsection*{Start fresh in the same folder}

Choose \key{1} (Start fresh session) and press \key{Enter} at
\prompt{Project name (.)} to keep working in the same folder. ProPrep starts a
new session file next to the old one. Nothing is loaded yet, so load the file
you want to continue from with the Structure Loader (\key{1a}, from a local
file; Chapter~\ref{ch:structure}).

\begin{watchout}
Starting ProPrep from the folder \emph{above} the project and typing the project
name again does not offer to resume. It opens the folder as a fresh project.
Always \typed{cd} into the project folder first.
\end{watchout}

Option \key{3} (Continue without session recording) works like a fresh start
but records nothing, so \key{undo} is not available.
